% =============================================================
% Practice Problems: Time Series Econometrics
% Course: Quantitative Economics, Vilnius University
% Author: Swapnil Singh
% =============================================================
\documentclass[10pt]{beamer}

\usepackage[T1]{fontenc}
\usepackage{graphicx}
\usepackage{booktabs}
\usepackage{tabularx}
\usepackage{array}
\usepackage{ragged2e}
\usepackage{appendixnumberbeamer}
\usepackage{amsmath,amssymb,amsfonts}
\usepackage{mathtools}
\usepackage{hyperref}
\usepackage{enumerate}

\usetheme{SimplePlus}
\graphicspath{{figures/time-series/}}

\title{Practice Problems: Time Series Econometrics}
\subtitle{Persistence, Trends, Forecasting, and Dynamic Systems}
\author{Swapnil Singh}
\institute{Lietuvos Bankas \and Kaunas University of Technology}
\date{Vilnius University, 2026}

\DeclareMathOperator{\Cov}{Cov}
\DeclareMathOperator{\Var}{Var}
\DeclareMathOperator{\Corr}{Corr}
\newcommand{\E}{\mathbb{E}}
\newcommand{\eps}{\varepsilon}
\newcommand{\calF}{\mathcal{F}}
\newcommand{\plim}{\operatorname*{plim}}
\newcommand{\I}{\mathbb{I}}

\begin{document}

\begin{frame}[plain]
  \titlepage
  \vfill
  {\scriptsize\color{gray} 8 problems from easy to hard. Based on \texttt{lectures/lecture\_time\_series.tex}.}
\end{frame}

\begin{frame}{How to use this practice set}
\small
\begin{block}{Structure}
The questions become progressively harder:
\begin{itemize}
  \item Problems 1--3: time-series language, stationarity, and autocorrelation.
  \item Problems 4--6: AR(1) mechanics, forecasting, and unit roots.
  \item Problems 7--8: time-series regression, serial correlation, and distributed lags.
\end{itemize}
\end{block}

\begin{alertblock}{Instruction}
Attempt each problem before opening the solution. The goal is not only to compute, but to state the maintained assumptions clearly.
\end{alertblock}
\end{frame}

% ============================================================
\section{Questions}
% ============================================================

\begin{frame}{Problem 1: Basic time-series language}
\label{q1}
\small
Let $\{Y_t:t=1,\ldots,T\}$ be a monthly time series for inflation.

\begin{enumerate}[(a)]
  \item What is the difference between a stochastic process and one observed time-series realization?
  \item Define the lag operator $L$ and write $\Delta Y_t$ using $L$.
  \item If $P_t$ is a price index, why is $\Delta \log P_t$ often interpreted as an inflation rate?
\end{enumerate}

\vfill
\hfill\hyperlink{s1}{\beamerbutton{Solution}}
\end{frame}

\begin{frame}{Problem 2: Weak stationarity}
\label{q2}
\small
Consider the process
\[
Y_t=\mu+\eps_t+0.5\eps_{t-1},
\]
where $\{\eps_t\}$ is white noise with $\E(\eps_t)=0$, $\Var(\eps_t)=\sigma^2$, and $\Cov(\eps_t,\eps_s)=0$ for $t\neq s$.

\begin{enumerate}[(a)]
  \item Compute $\E(Y_t)$.
  \item Compute $\Var(Y_t)$.
  \item Compute $\Cov(Y_t,Y_{t-1})$ and $\Cov(Y_t,Y_{t-2})$.
  \item Is $Y_t$ weakly stationary? Explain.
\end{enumerate}

\vfill
\hfill\hyperlink{s2}{\beamerbutton{Solution}}
\end{frame}

\begin{frame}{Problem 3: Reading autocorrelations}
\label{q3}
\small
Suppose a researcher computes sample autocorrelations for three variables:

\begin{center}
\begin{tabular}{lccc}
\toprule
Series & $\hat\rho_1$ & $\hat\rho_2$ & $\hat\rho_3$ \\
\midrule
A & 0.04 & -0.02 & 0.01 \\
B & 0.78 & 0.60 & 0.47 \\
C & 0.98 & 0.96 & 0.94 \\
\bottomrule
\end{tabular}
\end{center}

\begin{enumerate}[(a)]
  \item Which series looks closest to white noise?
  \item Which series looks persistent but plausibly stationary?
  \item Which series raises concern about a unit root or strong trend?
  \item Why is an ACF pattern only a diagnostic, not a proof?
\end{enumerate}

\vfill
\hfill\hyperlink{s3}{\beamerbutton{Solution}}
\end{frame}

\begin{frame}{Problem 4: AR(1) mean, variance, and ACF}
\label{q4}
\small
Let
\[
Y_t=c+\phi Y_{t-1}+\eps_t,\qquad \eps_t\sim WN(0,\sigma^2),
\]
with $|\phi|<1$.

\begin{enumerate}[(a)]
  \item Derive the unconditional mean $\mu=\E(Y_t)$.
  \item Derive the unconditional variance $\Var(Y_t)$.
  \item Derive $\rho_h=\Corr(Y_t,Y_{t-h})$.
  \item What happens to persistence as $\phi$ approaches one from below?
\end{enumerate}

\vfill
\hfill\hyperlink{s4}{\beamerbutton{Solution}}
\end{frame}

\begin{frame}{Problem 5: AR(1) forecasting}
\label{q5}
\small
Suppose quarterly inflation follows
\[
Y_t=1+0.6Y_{t-1}+\eps_t,
\qquad \E(\eps_{t+h}\mid\calF_t)=0.
\]
At date $T$, the observed inflation rate is $Y_T=4$.

\begin{enumerate}[(a)]
  \item Compute the unconditional mean of $Y_t$.
  \item Compute the one-step-ahead forecast $\widehat Y_{T+1|T}$.
  \item Compute the two-step-ahead forecast $\widehat Y_{T+2|T}$.
  \item Explain why the forecast converges to the unconditional mean as the horizon grows.
\end{enumerate}

\vfill
\hfill\hyperlink{s5}{\beamerbutton{Solution}}
\end{frame}

\begin{frame}{Problem 6: Unit root versus trend stationarity}
\label{q6}
\small
Consider two models:
\[
\text{A: } Y_t=\alpha+\delta t+u_t,\qquad u_t=\rho u_{t-1}+\eps_t,\quad |\rho|<1,
\]
\[
\text{B: } Y_t=Y_{t-1}+g+\eps_t.
\]

\begin{enumerate}[(a)]
  \item Which model is trend-stationary and which is difference-stationary?
  \item For each model, what transformation would you use before fitting a stationary ARMA model?
  \item What is the long-run effect of a one-time shock $\eps_t$ in each model?
  \item Why can confusing A and B lead to misleading empirical conclusions?
\end{enumerate}

\vfill
\hfill\hyperlink{s6}{\beamerbutton{Solution}}
\end{frame}

\begin{frame}{Problem 7: Serial correlation in a regression}
\label{q7}
\small
Suppose the true time-series regression is
\[
Y_t=\beta_0+\beta_1X_t+u_t,
\qquad
u_t=\rho u_{t-1}+e_t,
\]
where $e_t$ is white noise and $|\rho|<1$.

\begin{enumerate}[(a)]
  \item If $\E(u_t\mid X_1,\ldots,X_T)=0$, is OLS for $\beta_1$ consistent?
  \item If $\rho\neq 0$, are the usual homoskedastic OLS standard errors generally valid?
  \item Name one standard way to conduct inference robust to serial correlation.
  \item Explain the difference between inconsistency of $\hat\beta_1$ and invalid standard errors.
\end{enumerate}

\vfill
\hfill\hyperlink{s7}{\beamerbutton{Solution}}
\end{frame}

\begin{frame}{Problem 8: Distributed lags and long-run multipliers}
\label{q8}
\small
Consider the autoregressive distributed lag model
\[
Y_t=\alpha+\lambda Y_{t-1}+\delta_0X_t+\delta_1X_{t-1}+u_t,
\qquad |\lambda|<1.
\]

\begin{enumerate}[(a)]
  \item What is the immediate effect of a one-unit increase in $X_t$ holding $Y_{t-1}$ fixed?
  \item Derive the long-run multiplier of a permanent one-unit increase in $X$.
  \item Compute the long-run multiplier when $\lambda=0.5$, $\delta_0=0.2$, and $\delta_1=0.1$.
  \item Why is the long-run multiplier different from $\delta_0$?
\end{enumerate}

\vfill
\hfill\hyperlink{s8}{\beamerbutton{Solution}}
\end{frame}

% ============================================================
\appendix
\section{Solutions}
% ============================================================

\begin{frame}{Solution 1: Basic time-series language}
\label{s1}
\small
\begin{exampleblock}{Answer}
\begin{enumerate}[(a)]
  \item A stochastic process is the probabilistic sequence $\{Y_t\}$ that could generate many possible histories. A realization is the single path we observe in the data.
  \item The lag operator satisfies $LY_t=Y_{t-1}$. Therefore
  \[
  \Delta Y_t=Y_t-Y_{t-1}=(1-L)Y_t.
  \]
  \item Since
  \[
  \Delta\log P_t=\log P_t-\log P_{t-1}
  =\log(P_t/P_{t-1}),
  \]
  it is approximately the proportional growth rate of the price index. For small changes, it is close to $(P_t-P_{t-1})/P_{t-1}$.
\end{enumerate}
\end{exampleblock}
\vfill
\hfill\hyperlink{q1}{\beamerbutton{Back}}
\end{frame}

\begin{frame}{Solution 2: Weak stationarity}
\label{s2}
\small
\begin{exampleblock}{Moments}
\[
\E(Y_t)=\mu.
\]
Because $\eps_t$ and $\eps_{t-1}$ are uncorrelated,
\[
\Var(Y_t)=\Var(\eps_t+0.5\eps_{t-1})
=\sigma^2+0.25\sigma^2=1.25\sigma^2.
\]
Next,
\[
\Cov(Y_t,Y_{t-1})
=\Cov(\eps_t+0.5\eps_{t-1},\eps_{t-1}+0.5\eps_{t-2})
=0.5\sigma^2.
\]
For lag two,
\[
\Cov(Y_t,Y_{t-2})=0.
\]
\end{exampleblock}

\begin{block}{Conclusion}
The mean is constant, the variance is constant, and autocovariances depend only on the lag. So this MA(1) process is weakly stationary.
\end{block}
\vfill
\hfill\hyperlink{q2}{\beamerbutton{Back}}
\end{frame}

\begin{frame}{Solution 3: Reading autocorrelations}
\label{s3}
\small
\begin{exampleblock}{Classification}
\begin{enumerate}[(a)]
  \item Series A looks closest to white noise because its autocorrelations are near zero.
  \item Series B looks persistent but plausibly stationary: the ACF is positive and decays.
  \item Series C raises concern about a unit root or strong trend: the ACF stays very close to one and decays slowly.
  \item The sample ACF is only a diagnostic because finite samples, structural breaks, deterministic trends, seasonality, and nonlinear dynamics can produce similar-looking patterns.
\end{enumerate}
\end{exampleblock}

\begin{alertblock}{Econometric discipline}
The ACF tells us what to investigate next; it is not a substitute for a maintained model, plots, transformations, and formal tests.
\end{alertblock}
\vfill
\hfill\hyperlink{q3}{\beamerbutton{Back}}
\end{frame}

\begin{frame}{Solution 4: AR(1) mean, variance, and ACF}
\label{s4}
\footnotesize
\begin{exampleblock}{Mean}
\[
\mu=c+\phi\mu
\quad\Rightarrow\quad
\mu=\frac{c}{1-\phi}.
\]
\end{exampleblock}

\begin{exampleblock}{Variance and autocorrelation}
\[
Y_t-\mu=\phi(Y_{t-1}-\mu)+\eps_t.
\]
\[
\gamma_0=\phi^2\gamma_0+\sigma^2
\quad\Rightarrow\quad
\gamma_0=\frac{\sigma^2}{1-\phi^2}.
\]
For $h\geq 0$,
\[
\gamma_h=\phi^h\gamma_0,
\qquad
\rho_h=\frac{\gamma_h}{\gamma_0}=\phi^h.
\]
\end{exampleblock}

\begin{block}{Interpretation}
As $\phi\to 1^-$, shocks die out slowly and $\Var(Y_t)$ becomes large.
\end{block}
\vfill
\hfill\hyperlink{q4}{\beamerbutton{Back}}
\end{frame}

\begin{frame}{Solution 5: AR(1) forecasting}
\label{s5}
\small
\begin{exampleblock}{Calculations}
The unconditional mean is
\[
\mu=\frac{1}{1-0.6}=2.5.
\]
The one-step forecast is
\[
\widehat Y_{T+1|T}=1+0.6Y_T=1+0.6(4)=3.4.
\]
The two-step forecast is
\[
\widehat Y_{T+2|T}=1+0.6\widehat Y_{T+1|T}
=1+0.6(3.4)=3.04.
\]
\end{exampleblock}

\begin{block}{Mean reversion}
In general,
\[
\widehat Y_{T+h|T}=\mu+0.6^h(Y_T-\mu).
\]
Since $0.6^h\to 0$, the forecast converges to $\mu=2.5$.
\end{block}
\vfill
\hfill\hyperlink{q5}{\beamerbutton{Back}}
\end{frame}

\begin{frame}{Solution 6: Unit root versus trend stationarity}
\label{s6}
\small
\begin{exampleblock}{Classification and transformations}
\begin{enumerate}[(a)]
  \item Model A is trend-stationary: after removing $\alpha+\delta t$, the residual $u_t$ is stationary. Model B is difference-stationary: first differences satisfy $\Delta Y_t=g+\eps_t$.
  \item For A, detrend by regressing on a constant and time trend, then model the residual. For B, first-difference the series.
  \item In A, a one-time shock affects $u_t$ but dies out at rate $\rho^h$. In B, a one-time shock permanently shifts the level of $Y_t$.
  \item If a unit-root process is incorrectly detrended, one may mistake permanent shocks for temporary cycles. If a trend-stationary process is incorrectly differenced, one may throw away useful long-run level information.
\end{enumerate}
\end{exampleblock}
\vfill
\hfill\hyperlink{q6}{\beamerbutton{Back}}
\end{frame}

\begin{frame}{Solution 7: Serial correlation in a regression}
\label{s7}
\small
\begin{exampleblock}{Answer}
\begin{enumerate}[(a)]
  \item Yes, under strict exogeneity $\E(u_t\mid X_1,\ldots,X_T)=0$ and standard regularity conditions, OLS is consistent even if $u_t$ is serially correlated.
  \item No. Serial correlation generally invalidates the usual OLS standard errors because those formulas treat regression errors as uncorrelated across time.
  \item Use heteroskedasticity-and-autocorrelation consistent standard errors, such as Newey--West, or explicitly model the serial correlation when the maintained assumptions justify it.
  \item Inconsistency means $\hat\beta_1$ does not converge to $\beta_1$. Invalid standard errors mean the estimator may be centered correctly, but the reported uncertainty, $t$ statistics, and confidence intervals are wrong.
\end{enumerate}
\end{exampleblock}
\vfill
\hfill\hyperlink{q7}{\beamerbutton{Back}}
\end{frame}

\begin{frame}{Solution 8: Distributed lags and long-run multipliers}
\label{s8}
\small
\begin{exampleblock}{Immediate and long-run effects}
\begin{enumerate}[(a)]
  \item Holding $Y_{t-1}$ fixed, the immediate effect of $X_t$ is $\delta_0$.
  \item In a long-run steady state, set $Y_t=Y_{t-1}=Y$ and $X_t=X_{t-1}=X$:
  \[
  Y=\alpha+\lambda Y+(\delta_0+\delta_1)X.
  \]
  Hence
  \[
  \frac{\partial Y}{\partial X}
  =\frac{\delta_0+\delta_1}{1-\lambda}.
  \]
  \item With $\lambda=0.5$, $\delta_0=0.2$, and $\delta_1=0.1$,
  \[
  \frac{0.2+0.1}{1-0.5}=0.6.
  \]
  \item The long-run multiplier includes the direct lag effect $\delta_1$ and the propagation through lagged $Y$ governed by $\lambda$.
\end{enumerate}
\end{exampleblock}
\vfill
\hfill\hyperlink{q8}{\beamerbutton{Back}}
\end{frame}

\end{document}
